% ─────────────────────────────────────────────────────────────────────────────
%  SECTION 3 — THE BAILOUT PROTOCOL
%  Formal specification: state machine, trigger conditions, escalation algorithm,
%  bypass mechanism, and timeout handling.
% ─────────────────────────────────────────────────────────────────────────────

\section{The Bailout Protocol}\label{sec:bailout}

% ─────────────────────────────────────────────────────────────────────────────
%  3.1  Overview and Architectural Purpose
% ─────────────────────────────────────────────────────────────────────────────

\subsection{Overview and Architectural Purpose}\label{sec:bailout-overview}

The Bailout Protocol is the mandatory exception-escalation mechanism embedded in
every node of a \bxthree{} system. It is not an error-handling routine selected
at runtime, nor an optional escalation pathway available to a system that chooses
to invoke it. It is an \emph{architectural compulsion}: any condition a node
cannot resolve with certainty must propagate to a Human Accountability Anchor,
and no machine actor in the propagation chain may absorb, resolve, or suppress
that propagation.

The protocol's purpose is threefold. First, it provides a deterministic state
machine that governs how an agent node transitions from nominal operation through
exception detection to human resolution. Second, it defines a bypass mechanism
that guarantees the propagation chain cannot be intercepted, delayed, or absorbed
by any machine actor—ensuring that human reachability is architecturally
guaranteed rather than policy-dependent. Third, it enforces a timeout hierarchy
that bounds the worst-case wall-clock time from trigger firing to resolution,
providing a formal liveness guarantee.

This section provides the complete formal specification: the state machine,
the precise transition conditions, the bail-out trigger condition, the escalation
algorithm, the bypass mechanism, and the timeout hierarchy.

% ─────────────────────────────────────────────────────────────────────────────
%  3.2  State Machine Formal Definition
% ─────────────────────────────────────────────────────────────────────────────

\subsection{State Machine Definition}\label{sec:bailout-state-machine}

The Bailout Protocol is formally modeled as a deterministic finite state machine
$\mathcal{S}_{\text{bailout}}$ embedded in the \fact{} layer of every \bxthree{}
node:

\[
\mathcal{S}_{\text{bailout}} \;=\; \bigl( Q,\; \Sigma,\; \delta,\; q_0,\; F \bigr)
\]

\begin{itemize}\itemsep=0pt
\item[$Q$] $\{\;\text{IDLE},\; \text{BOUNDED},\; \text{BAIL\_PENDING},\;
  \text{ESCALATING},\; \text{HUMAN\_ACKNOWLEDGED},\; \text{RESOLVED}\; \}$ is the
  finite set of protocol states.
\item[$\Sigma$] The input alphabet defined in Table~\ref{tab:bailout-events}.
\item[$\delta : Q \times \Sigma \to Q$] The deterministic transition function
  (Table~\ref{tab:bailout-transition}). For every $(q, \sigma) \in Q \times
  \Sigma$, at most one $q' \in Q$ satisfies $\delta(q, \sigma) = q'$.
  If no transition is defined for $(q, \sigma)$, the machine stalls and a
  protocol violation is recorded in the Forensic Ledger.
\item[$q_0$] $\text{IDLE}$ is the designated initial state, entered at node
  startup and after every successful directive resolution.
\item[$F$] $\{\;\text{RESOLVED}\; \}$ is the set of terminal accepting states.
\end{itemize}

A node's state $q \in Q$ is stored in the node's \fact{} layer worksheet and is
atomically observable by its parent node and by the Bailout Monitor. State updates
are atomic: no intermediate or partially-executing states are externally
observable. The machine is \emph{deterministic} by construction; the Bailout
Monitor's priority function $\pi$ totally orders simultaneous trigger conditions
and selects the highest-priority enabled transition before committing any state
change to the Forensic Ledger.

\bigskip

\noindent\textbf{State definitions:}

\begin{trivlist}\itemindent=2em
\item[\textbf{IDLE}] The node's Bounds Engine is operating within its provisioned
  capability. No exceptional condition exists. Normal sensing, modelling, and
  proposition cycles are active.

\item[\textbf{BOUNDED}] The node has detected a condition that \emph{may} require
  escalation, but the Bounds Engine is still evaluating whether the condition is
  resolvable locally or requires human intervention. The condition is known and
  contained. No propagation has occurred.

\item[\textbf{BAIL\_PENDING}] The node has determined the condition cannot be
  resolved locally. The trigger package $\mathcal{T}$ has been assembled and is
  in the process of being propagated to the parent node. Delivery is asynchronous
  and confirmed.

\item[\textbf{ESCALATING}] The trigger package has been delivered to a machine
  actor in the propagation chain. That actor is evaluating whether it can resolve
  the exception. If it cannot, it propagates further upward. This state persists
  across zero or more machine-actor hops.

\item[\textbf{HUMAN\_ACKNOWLEDGED}] The trigger package has reached a Human
  Accountability Anchor. The human operator has received the alert and is actively
  reviewing the exception. The system awaits a directive from the human.

\item[\textbf{RESOLVED}] The human has issued a directive and the directive has
  propagated back to the originating node. The actuators have executed or will
  execute the directive. The trigger is closed and the complete lifecycle is
  sealed in the Forensic Ledger.
\end{trivlist}

% ─────────────────────────────────────────────────────────────────────────────
%  3.3  Transition Conditions and Transition Table
% ─────────────────────────────────────────────────────────────────────────────

\subsection{Transition Conditions}\label{sec:bailout-transitions}

Let $e \in \Sigma$ denote the event driving a transition. Table
\ref{tab:bailout-transition} defines $\delta(q, e)$. Entries marked ``—''
are undefined transitions: if a machine actor in state $q$ receives event $e$
with no defined transition, the machine stalls and the Bailout Monitor logs
a $\texttt{TRANSITION\_STALL}$ incident.

\begin{table}[H]
\centering
\caption{Bailout Protocol Transition Function $\delta(q, e)$}
\label{tab:bailout-transition}
\begin{tabular}{lcccccc}
\toprule
Event $e \backslash$ State $q$ & IDLE & BOUNDED & BAIL\_PENDING &
  ESCALATING & HUMAN\_ACK & RESOLVED \\
\midrule
\texttt{dispatched}           & BOUNDED & —    & —    & —    & —    & — \\
\texttt{tc\_fire}            & —       & BOUNDED & BAIL\_PENDING & — & — & — \\
\texttt{bounds\_exceeded}    & —       & BAIL\_PENDING & — & — & — & — \\
\texttt{delivery\_confirm}    & —       & —    & ESCALATING & — & — & — \\
\texttt{hop\_evaluated}       & —       & —    & —    & ESCALATING & — & — \\
\texttt{propagate\_up}       & —       & —    & —    & ESCALATING & — & — \\
\texttt{human\_received}     & —       & —    & —    & —    & HUMAN\_ACK & — \\
\texttt{directive\_issued}   & —       & —    & —    & —    & RESOLVED & — \\
\texttt{timeout\_escape}     & —       & —    & ESCALATING & — & — & — \\
\texttt{timeout\_human}      & —       & —    & —    & RESOLVED & — & — \\
\texttt{resolve\_local}       & —       & IDLE & —    & —    & — & — \\
\texttt{directive\_complete} & —       & —    & —    & —    & — & RESOLVED \\
\bottomrule
\end{tabular}
\end{table}

Each event is formally defined as follows:

\begin{trivlist}\itemindent=0em
\item[\texttt{dispatched}:] The node has been activated by the scheduler.
  $\text{dispatched}(n) = \text{true}$.

\item[\texttt{tc\_fire}:] The trigger condition TC has evaluated to true at node
  $n$. See Section~\ref{sec:bailout-trigger}.

\item[\texttt{bounds\_exceeded}:] The Bounds Engine has determined that the
  condition at node $n$ cannot be resolved within the node's Safety Envelope.
  Formally: $\text{BoundsEngine.evaluate}(n, \mathcal{T}).\text{within\_SE} =
  \text{false}$.

\item[\texttt{delivery\_confirm}:] The trigger package $\mathcal{T}$ has been
  successfully delivered to and acknowledged by the parent node $\alpha(n)$.

\item[\texttt{hop\_evaluated}:] A machine actor in the propagation chain has
  completed evaluation of $\mathcal{T}$ and found it cannot resolve the
  exception locally.

\item[\texttt{propagate\_up}:] The evaluating machine actor propagates $\mathcal{T}$
  to its own parent $\alpha(\text{actor})$.

\item[\texttt{human\_received}:] The trigger package has been delivered to a
  Human Accountability Anchor $h(n)$ and the human has acknowledged receipt.
  $\text{isHuman}(h(n)) = \text{true} \land \text{ack}(h(n), \mathcal{T}) =
  \text{true}$.

\item[\texttt{directive\_issued}:] The human operator has issued a directive $d$
  to resolve the exception. $d \in \mathcal{D}$, where $\mathcal{D}$ is the set
  of permissible directives defined in the protocol specification.

\item[\texttt{timeout\_escape}:] The elapsed time $\tau_{\text{escape}}(n)$
  since entering \texttt{BAIL\_PENDING} has exceeded $T_{\text{escape}}$.
  $\tau_{\text{escape}}(n) > T_{\text{escape}}$.

\item[\texttt{timeout\_human}:] The elapsed time $\tau_{\text{human}}(n)$ since
  entering \texttt{HUMAN\_ACKNOWLEDGED} has exceeded $T_{\text{human}}$.
  $\tau_{\text{human}}(n) > T_{\text{human}}$.

\item[\texttt{resolve\_local}:] The Bounds Engine has produced a resolution
  within the Safety Envelope. The node can continue without human intervention.
  $\text{BoundsEngine.resolve}(n, \mathcal{T}).\text{within\_SE} = \text{true}$.

\item[\texttt{directive\_complete}:] The originating node has executed the
  human directive and the actuators have reached the commanded state.
\end{trivlist}

% ─────────────────────────────────────────────────────────────────────────────
%  3.4  Bail-Out Trigger Condition
% ─────────────────────────────────────────────────────────────────────────────

\subsection{Bail-Out Trigger Condition}\label{sec:bailout-trigger}

The bail-out trigger condition TC is the sole entry criterion for the
BOUNDED $\rightarrow$ BAIL\_PENDING transition. It is evaluated continuously
by the \bounds{} layer at each scheduling tick and is defined as the disjunction
of four independent sub-conditions:

\begin{equation}
\text{TC}(n) \;=\; \rho_{\text{budget}}(n) \;\lor\; \rho_{\text{time}}(n)
\;\lor\; \rho_{\text{safety}}(n) \;\lor\; \rho_{\text{accountability}}(n)
\tag{eq:trigger-condition}
\end{equation}

\begin{trivlist}\itemindent=0em
\item[\textbf{Resource budget exhaustion $\rho_{\text{budget}}$:}]
\[
\rho_{\text{budget}}(n) \;\iff\; \text{consumed}(n) \;>\; \text{budget}(n)
\]
The cumulative resource consumption of node $n$ exceeds its allocated budget.
This includes compute units, memory allocations, and I/O operations tracked by
the \fact{} layer's resource monitor.

\item[\textbf{Execution time limit $\rho_{\text{time}}$:}]
\[
\rho_{\text{time}}(n) \;\iff\; \tau_{\text{exec}}(n) \;>\; T_{\text{bounds}}(n)
\]
The cumulative execution time $\tau_{\text{exec}}$ since node activation exceeds
the configured time bound $T_{\text{bounds}}$. This bound is node-type-specific
and is immutable for safety-critical node types.

\item[\textbf{Safety envelope prediction $\rho_{\text{safety}}$:}]
\[
\rho_{\text{safety}}(n) \;\iff\; \Pr\!\bigl(\text{action}(n) \notin \text{SafetyEnvelope}\bigr) \;>\; \theta_{\text{safety}}
\]
The probabilistic safety predictor embedded in the \bounds{} layer assigns a
probability greater than $\theta_{\text{safety}}$ (default: $0.05$) to the next
proposed action exiting the node's Safety Envelope. This is a predictive trigger:
it fires before the action executes, preventing unsafe state transitions.

\item[\textbf{Accountability boundary violation $\rho_{\text{accountability}}$:}]
\[
\rho_{\text{accountability}}(n) \;\iff\; \text{objects\_accessed}(n) \;\not\subseteq\; \text{objects\_declared}(n)
\]
The node has accessed or modified objects outside its declared capability scope.
This covers scope creep, capability drift, and tooling injection attacks. The
condition is detected by the \fact{} layer's audit hook at each object access.
\end{trivlist}

TC is evaluated \emph{deterministically} and \emph{side-effect-free}: the
trigger result is written to the task's runtime state and broadcast to the
Bailout Monitor via the event bus. The trigger cannot be suppressed, delayed,
or rewritten by any machine actor including the node itself. The \fact{} layer
pre-commit hook enforces this: any attempt to modify the trigger evaluation
routine or its result is logged as a $\texttt{TRIGGER\_TAMPERING}$ security
event in the Forensic Ledger.

% ─────────────────────────────────────────────────────────────────────────────
%  3.5  Escalation Algorithm
% ─────────────────────────────────────────────────────────────────────────────

\subsection{Escalation Algorithm}\label{sec:bailout-escalation}

When TC fires and the node transitions to BOUNDED, the Bounds Engine has
$T_{\text{escape}}$ (default: 300 seconds) to either resolve the condition
locally within the Safety Envelope or determine that it cannot. If the
condition is unresolved at timeout, the node executes the following escalation
algorithm to route the exception to the Human Accountability Anchor. The
algorithm is embedded identically in every node's Worksheet; it is
deterministic and cannot be overridden by any machine actor.

\begin{algorithm}[H]
\caption{Bailout Protocol Escalation Algorithm (executed at every node)}
\label{alg:bailout-escalation}
\begin{enumerate}
\item \textbf{Initialize.} Set $\text{current} \leftarrow n$,
  $\text{propagation\_chain} \leftarrow []$, $r \leftarrow 0$,
  $\text{hop\_count} \leftarrow 0$.

\item \textbf{Check terminus.} If $\text{isHuman}(\text{current}) = \text{true}$,
  go to Step~\ref{step:deliver}.

\item \textbf{Record propagation step.} Append
  $\{\text{node}: \text{current.id},\; \text{action}: \texttt{ESCALATING},\;
  t_{\text{now}},\; \text{hop}: \text{hop\_count}\}$ to
  $\text{propagation\_chain}$.

\item \textbf{Evaluate at machine actor.} Set
  $r \leftarrow \text{current.BoundsEngine.evaluate}(\mathcal{T})$.
  \begin{enumerate}
  \item If $r.\text{can\_resolve} = \text{true}$ and
    $r.\text{within\_SafetyEnvelope} = \text{true}$:
    \begin{enumerate}
    \item $\mathcal{T}.\text{resolution} \leftarrow r$
    \item $\mathcal{T}.\text{status} \leftarrow \texttt{RESOLVED\_LOCAL}$
    \item Propagate resolution back down the chain to $n$; \textbf{return}
    \end{enumerate}
  \item Else: continue to Step~\ref{step:parent}.
  \end{enumerate}

\item \label{step:parent}\textbf{Move to parent.} Set $\text{current}
  \leftarrow \alpha(\text{current})$ (the immutable parent pointer).
  Increment $\text{hop\_count}$.

\item \textbf{Hard hop limit.} If $\text{hop\_count} > \text{HOP\_LIMIT}$
  (default: 127), transition the originating node to
  $\texttt{QUIESCENT}$ and alert the Human Root via fallback channels;
  record $\texttt{ESCALATION\_LOOP\_PREVENTED}$ in the Forensic Ledger;
  \textbf{return}.

\item \textbf{Retry with backoff.} Set $\text{delay} \leftarrow
  \text{Backoff}(\text{hop\_count}, T_{\text{prop}})$; wait \text{delay};
  attempt delivery to $\text{current}$. If delivery fails, go to
  Step~\ref{step:parent}.

\item \label{step:deliver}\textbf{Deliver to human.} Set
  $\mathcal{T}.\text{status} \leftarrow \texttt{HUMAN\_ACKNOWLEDGED}$.
  Deliver $\mathcal{T}$ to the human interface (Root Dashboard alert,
  SMS, or fallback). Await directive $d$ from the human operator.

\item \textbf{Process human response.} If $d$ received before
  $T_{\text{human}}$:
  \begin{enumerate}
  \item $\mathcal{T}.\text{directive} \leftarrow d$
  \item $\mathcal{T}.\text{status} \leftarrow \texttt{DIRECTIVE\_ISSUED}$
  \item Propagate $d$ back down $\text{propagation\_chain}$ to $n$
  \item Transition $n$ to \texttt{RESOLVED}; \textbf{return}
  \end{enumerate}
  Else: transition to \texttt{RESOLVED} with tag
  $\texttt{RESOLVED\_TIMEOUT}$; record in Forensic Ledger; \textbf{return}.
\end{enumerate}
\end{algorithm}

\bigskip
\noindent\textbf{Backoff function.} The retry backoff follows an exponential
schedule to prevent thundering-herd load on parent nodes during network
partition:

\begin{equation}
\text{Backoff}(h, T_{\text{prop}}) \;=\; \min\!\bigl(
  T_{\text{prop}} \cdot 2^{h},\;
  T_{\text{cap}}
\bigr)
\tag{eq:backoff}
\end{equation}

Where $h$ is the current hop count and $T_{\text{cap}}$ is the per-step
cap (default: 30 seconds). The backoff is computed by the Bailout Monitor,
not by any machine actor, preventing artificial backoff manipulation.

\bigskip
\noindent\textbf{Propagation chain integrity.} The propagation chain
$\text{propagation\_chain}$ is a timestamped, ordered list of every node
that received and evaluated $\mathcal{T}$. It is written incrementally to the
Forensic Ledger at each hop, enabling post-hoc reconstruction of the full
escalation path. The chain is immutable once written: no node may edit,
truncate, or delete its entries.

% ─────────────────────────────────────────────────────────────────────────────
%  3.6  Bypass Mechanism
% ─────────────────────────────────────────────────────────────────────────────

\subsection{Bypass Mechanism}\label{sec:bailout-bypass}

The bypass mechanism is the architectural guarantee that the Bailout Protocol
reaches a human even when machine actors in the propagation chain are
malfunctioning, adversarial, or compromised. It exploits two immutable
structures embedded in every node's worksheet: the immutable parent pointer
$\alpha(n)$ and the Human Accountability Anchor reference $h(n)$.

\bigskip
\noindent\textbf{Why it bypasses all machine actors:}

The bypass mechanism defeats three specific failure modes that could otherwise
prevent escalation:

\begin{trivlist}\itemindent=0em
\item[\textbf{Intercept:}] A machine actor in the propagation chain receives
  $\mathcal{T}$ and deliberately does not propagate it, instead issuing a
  false $\texttt{RESOLVED\_LOCAL}$ tag. The bypass mechanism defeats this because
  the machine actor cannot suppress the immutable log entry: the Bailout Monitor
  records every hop. If a machine actor fails to propagate within
  $T_{\text{prop}}$ of receiving $\mathcal{T}$, the Monitor itself escalates to
  the next parent, bypassing the non-responsive actor entirely.

\item[\textbf{Redirect:}] A compromised parent node $\alpha(n)$ attempts to
  redirect $\mathcal{T}$ to a decoy human interface under attacker control. The
  bypass mechanism defeats this because each node's $h(n)$ reference is sealed
  in the \fact{} layer at node creation and cannot be rewritten by any machine
  actor. The propagation chain follows the sealed reference, not a routing table
  an attacker could manipulate.

\item[\textbf{Delay:}] A machine actor propagates $\mathcal{T}$ but introduces
  an unbounded delay before forwarding it to the parent, effectively suppressing
  escalation by making it timing-unreliable. The bypass mechanism defeats this
  because the propagation timeout $T_{\text{prop}}$ is enforced by the Bailout
  Monitor, not by the machine actor. If delivery is not confirmed within
  $T_{\text{prop}}$, the Monitor independently retries to the parent, bypassing
  the delayed actor. The backoff curve (Equation~\ref{eq:backoff}) prevents
  thundering-herd load while still guaranteeing eventual delivery.
\end{trivlist}

\bigskip
\noindent\textbf{Mechanism description:}

The bypass is implemented as a separate Bailout Monitor thread that operates
independently of the node's main execution context. The Monitor maintains:

\begin{enumerate}
\itemsep=0pt
\item A sealed copy of $h(n)$ for every descendant node, propagated downward
  at node creation and stored in the \fact{} layer.
\item An independent communication channel to each $h(n)$, provisioned at
  system initialization and independent of the node's primary messaging stack.
  If the primary stack is corrupted (e.g., by a compromised node), the Monitor
  uses the fallback channel.
\item A propagation timer for each active $\mathcal{T}$. If the timer expires
  before $\texttt{delivery\_confirm}$ is recorded, the Monitor escalates directly
  to $\alpha(n)$, bypassing any machine actor that has failed to propagate.
\end{enumerate}

The Machine Actor Exclusion Invariant (MAE) formally guarantees:

\begin{center}
\fbox{
\begin{minipage}{0.92\linewidth}
\textbf{Invariant MAE:} For any node $n$ and any trigger package
$\mathcal{T}$ generated at $n$, there exists no machine actor $m$ in the
propagation chain from $n$ to $h(n)$ such that $m$ can prevent the delivery of
$\mathcal{T}$ to $h(n)$ within $T_{\max}$ wall-clock time from TC firing.
\end{minipage}
}
\end{center}

MAE is enforced by architectural constraints, not by policy or configuration:
the Bailout Monitor thread runs in the \fact{} layer with higher privilege than
any machine actor, and its propagation logic is defined in the protocol
specification (Algorithm~\ref{alg:bailout-escalation}) with no conditional
branches that could be redirected by a machine actor.

% ─────────────────────────────────────────────────────────────────────────────
%  3.7  Timeout Handling
% ─────────────────────────────────────────────────────────────────────────────

\subsection{Timeout Handling}\label{sec:bailout-timeouts}

Timeouts are the backbone of the Bailout Protocol's liveness guarantees. Without
timeouts, a node in \texttt{BAIL\_PENDING} could remain in that state
indefinitely, and an escalation chain could propagate indefinitely without ever
reaching a human. The protocol defines a four-level timeout hierarchy that bounds
the worst-case wall-clock time from trigger firing to human acknowledgment:

\bigskip
\noindent\textbf{Level 1 — Bounds evaluation timeout $T_{\text{bounds}}$:}
After TC fires and the node enters \texttt{BOUNDED}, the Bounds Engine has at
most $T_{\text{bounds}}$ (default: 300 seconds) to either resolve the condition
locally or confirm it cannot. If the condition persists at $T_{\text{bounds}}$,
the transition to \texttt{BAIL\_PENDING} is forced by the \fact{} pre-commit
hook regardless of the Bounds Engine's internal state. This timeout cannot be
extended by any machine actor.

\bigskip
\noindent\textbf{Level 2 — Propagation timeout $T_{\text{prop}}$:}
At each machine-actor hop in the escalation chain, the Bailout Monitor allows
$T_{\text{prop}}$ (default: 10 seconds, with exponential backoff on retry) for
delivery confirmation. If delivery fails, the Monitor independently retries
to the same parent up to $\text{MAX\_RETRY}$ times (default: 3) before moving
to the next parent in the chain. The backoff curve (Equation~\ref{eq:backoff})
prevents network saturation while bounding total propagation time.

\bigskip
\noindent\textbf{Level 3 — Pathway exhaustion cap $T_{\text{cap}}$:}
If all communication pathways to the parent fail consecutively for
$\text{MAX\_FAIL} \cdot T_{\text{prop}}$ (default: 90 seconds), the protocol
marks the pathway as exhausted and transitions the originating node to
$\texttt{QUIESCENT}$—a suspended state with all actuators paused—while
alerting the Human Root via the fallback channel. The tag
$\texttt{PATHWAY\_EXHAUSTED}$ is recorded in the Forensic Ledger. This prevents
indefinite retry loops on partitioned networks.

\bigskip
\noindent\textbf{Level 4 — Human acknowledgment timeout $T_{\text{human}}$:}
After $\mathcal{T}$ reaches the Human Accountability Anchor and the node enters
$\texttt{HUMAN\_ACKNOWLEDGED}$, the system awaits a directive for a maximum of
$T_{\text{human}}$ (default: 3600 seconds). If no directive is received:

\begin{enumerate}\itemsep=0pt
\item The system transitions to $\texttt{RESOLVED}$ with tag
  $\texttt{RESOLVED\_TIMEOUT}$.
\item The Human Root is alerted via all available fallback pathways.
\item The termination reason is immutably recorded in the Forensic Ledger,
  ensuring the failure is attributed as a human non-response incident requiring
  post-mortem review—not a silent suppression.
\end{enumerate}

\bigskip
\noindent\textbf{Worst-case liveness bound:} Together, the four timeouts define
a finite upper bound $T_{\max}$ from TC firing to human acknowledgment under
any combination of parent unavailability and pathway degradation:

\begin{equation}
T_{\max} \;=\; T_{\text{bounds}}
\;+\; \text{HOP\_LIMIT} \cdot T_{\text{prop}} \cdot 2^{\text{HOP\_LIMIT}}
\;+\; T_{\text{human}}
\tag{eq:tmax}
\end{equation}

A deployment that cannot guarantee human acknowledgment within $T_{\max}$ must
provision additional or higher-bandwidth pathways, or must reduce
$T_{\text{prop}}$ to bring the worst-case bound within the required limit.
This analysis is a required part of deployment certification for any \bxthree{}
system operating in a regulated or safety-critical environment.

\bigskip
\noindent\textbf{Timeout interaction with the bypass mechanism:} The bypass
mechanism ensures that timeouts cannot be weaponised by a machine actor to
delay or suppress escalation:

\begin{itemize}\itemsep=0pt
\item $T_{\text{bounds}}$ is enforced by the \fact{} layer pre-commit hook,
  not by the \bounds{} layer. A machine actor cannot artificially extend
  $T_{\text{bounds}}$ to prevent transition to \texttt{ESCALATING}.
\item $T_{\text{prop}}$ backoff is calculated by the Bailout Monitor, not
  by any machine actor. The backoff curve is defined in the protocol
  specification and cannot be altered at runtime.
\item If all pathways timeout at $T_{\text{cap}}$, the protocol resolves with
  $\texttt{PATHWAY\_EXHAUSTED}$—a terminal \texttt{RESOLVED} state, not a
  suppression of escalation. The tag in the Forensic Ledger ensures the
  failure is recorded as an incident, not silent suppression.
\end{itemize}

% ─────────────────────────────────────────────────────────────────────────────
%  3.8  Summary of Protocol Properties
% ─────────────────────────────────────────────────────────────────────────────

\subsection{Summary of Protocol Properties}\label{sec:bailout-summary}

The Bailout Protocol satisfies the following formal properties:

\begin{enumerate}\itemsep=0pt
\item $\mathcal{S}_{\text{bailout}}$ is a deterministic finite state machine
  with six states and a single terminal accepting state (\texttt{RESOLVED}).
  Every state has at most one defined outbound transition per event.

\item The trigger condition TC (Equation~\ref{eq:trigger-condition}) is
  individually necessary and jointly sufficient to cover all conditions
  requiring human escalation: resource budget exhaustion, execution time
  limit, safety envelope prediction, and accountability boundary violation.

\item The Machine Actor Exclusion Invariant (MAE) is enforced by the
  transition function $\delta$ and the Bailout Monitor's independent
  operation. It cannot be disabled, relaxed, or circumvented by any
  machine actor. The bypass mechanism eliminates the Intercept, Redirect,
  and Delay failure modes by architectural constraint.

\item The escalation algorithm (Algorithm~\ref{alg:bailout-escalation})
  terminates within $\text{distance}(n, h(n)) + 1$ propagation hops, or
  triggers $\texttt{QUIESCENT}$ halt on pathway exhaustion. No infinite
  loops are possible in the propagation chain.

\item Every state transition is immutably recorded in the Forensic Ledger
  with a SHA-256 digest chain providing tamper-evident auditability.

\item The four-level timeout hierarchy
  ($T_{\text{bounds}}$, $T_{\text{prop}}$, $T_{\text{cap}}$,
  $T_{\text{human}}$) guarantees a finite worst-case liveness bound
  $T_{\max}$ (Equation~\ref{eq:tmax}) for every protocol run.

\item The bypass mechanism guarantees that human reachability survives
  delegation depth, network partition, and recursive node spawning, by
  exploiting the immutable parent-pointer structure and Human Accountability
  Anchor reference embedded in each node's worksheet.
\end{enumerate}

These properties collectively ensure that the Bailout Protocol delivers on
its architectural guarantee: \emph{any} condition a machine actor cannot
resolve with certainty \emph{must} propagate to a Human Accountability
Anchor, and no machine actor in the propagation chain may suppress,
absorb, redirect, or delay that propagation. The Forensic Ledger ensures
that every run—successful or timed-out—is permanently, immutably recorded.